\section{Seguimiento operativo y notas internas}

Esta sección complementa la lectura de resultados con un registro práctico de desarrollo: qué se ha generado, qué queda pendiente y qué cuestiones metodológicas siguen abiertas.

\subsection{Estado de los artefactos generados}

El pipeline actual emite salidas de control de calidad y análisis en cada etapa: diagnósticos de ingestión y filtrado, resúmenes de componentes ICA y tablas de evaluación, estadísticas de segmentación/baseline/rechazo de artefactos, visualizaciones y matrices FFT/CSD/wPLI, y tablas de inferencia basada en surrogates. Mantener estos outputs agrupados por fase es esencial para depurar con rapidez, porque permite validar supuestos locales antes de propagar cambios a etapas posteriores.
\sidenote{%
\footnotesize
\textbf{Mínimo por ejecución}\\
1) logs por fase\\
2) figuras QC\\
3) tablas resumen\\
4) snapshot de config
}

\begin{fullwidth}
\begin{figure}[H]
\centering
\resizebox{0.74\fulltextwidth}{!}{\NSOutputsDashboard}
\caption{Dashboard conceptual de outputs por ejecución. Una ejecución de referencia debe poder reconstruirse con logs, figuras QC, tablas resumen, snapshot de configuración, hash/commit y carpeta de resultados inequívoca.}
\label{fig:generated-outputs-dashboard}
\end{figure}
\end{fullwidth}

\subsection{Tareas pendientes}

Los elementos pendientes son principalmente operativos: verificar completitud y consistencia de etiquetado EO/EC en todos los participantes, fijar el número final de surrogates y relanzar la inferencia con configuración congelada, y consolidar las salidas por banda en tablas internas reproducibles de comparación. Estas tareas deben registrarse con timestamp de ejecución y hash/configuración para evitar ambigüedades entre ejecuciones exploratorias y ejecuciones de referencia.
\sidenote{%
\footnotesize
\textbf{Etiqueta recomendada}\\
\texttt{run\_<date>\_<branch>\_<cfg>}\\
Incluir commit hash y parámetros clave.
}

\subsection{Hipótesis de trabajo y preguntas abiertas}

Las hipótesis actuales se centran en la modulación esperada EO/EC de la conectividad relacionada con alfa, la mejora de especificidad esperada al activar CSD antes de wPLI frente a alternativas sin CSD, y la estabilización de la densidad de aristas significativas cuando se aplican umbrales por surrogates. Las preguntas técnicas abiertas incluyen la sensibilidad a la longitud/solapamiento de segmentos, el posible exceso de limpieza en sujetos de baja SNR durante el rechazo de artefactos por ICA, y el compromiso entre tiempo de ejecución y precisión al aumentar el número de surrogates.
